% Suggested replacement for the Related Work section; primary-source audit:
% docs/paper/writing/sources.md. Approximately 580 prose words.
\section{Related Work}
\label{sec:related-work}

\subsection{Slicing, Feature Location, and Software Reuse}

Feature lifting builds on longstanding research into recovering reusable behavior
from existing software. Program slicing uses data and control dependencies to
retain behavior with respect to a specified slicing
criterion~\cite{weiser1984slicing}. Feature-location techniques identify the source
entities associated with a requested capability; the taxonomy of Dit et~al.
organizes the evidence and techniques used for this
task~\cite{dit2013featurelocation}. Multi-dimensional separation of concerns
addresses decompositions in which concerns overlap rather than align with a
single module hierarchy~\cite{tarr1999degrees}. These perspectives explain why a
feature can cross file and package boundaries and why locating an entry point
does not by itself establish a complete extraction.

Automated software transplantation is a particularly close antecedent. Barr
et~al. combine program analysis, lightweight annotations, testing, and search to
extract functionality from a donor and adapt it for a different
host~\cite{barr2015transplantation}. Their work already makes preservation during
extraction and integration an executable objective. \flb carries this reuse
objective into the evaluation of general-purpose repository agents: the agent
receives a behavioral contract and an intact donor repository, determines the
supporting implementation without source-location hints, and delivers a standalone
package. It does not receive a destination application's integration site. We
evaluate finite behavioral obligations and source independence, without claiming
to establish semantic equivalence or compute a minimal program slice.

\subsection{Executable Coding and Feature-Development Benchmarks}

HumanEval measures functional correctness for code synthesized from
docstrings~\cite{chen2021codex}. SWE-bench moves evaluation into real repositories,
requiring changes that resolve existing GitHub
issues~\cite{jimenez2024swebench}. More recent benchmarks directly address feature
development. FeatBench uses natural-language requirements without code hints to
evaluate feature implementation in existing
repositories~\cite{chen2026featbench}. FeatureBench derives tasks through
test-driven dependency tracing and evaluates both incremental implementation in a
repository and construction of functionality from
scratch~\cite{zhou2026featurebench}. Thus, producing a callable module or
implementing a repository-derived feature is already represented in prior
evaluation settings.

The distinguishing experimental condition in \flb is access to the complete
existing implementation as evidence, coupled with a requirement to preserve its
declared behavior across a new package boundary. The final artifact must remain
functional when access to the donor repository is removed. This condition makes
decisions about transitive helpers, resources, state, and runtime dependencies
part of the evaluated behavior. Our Full-Repository / No-Hint protocol withholds
source-location hints while specifying the required output API, so it should not
be equated with withholding all interface information. The benchmark additionally
reports extraction footprint on successful artifacts; reference-relative size is
a descriptive measure, not a certificate of minimality or maintainability.

\subsection{Agent Runtimes and Evaluation Scope}

SWE-agent demonstrates that the agent--computer interface affects how effectively
language models navigate repositories, edit code, and execute
tests~\cite{yang2024sweagent}. OpenHands provides an extensible platform for agents
that use code execution, command-line interaction, and other tools in sandboxed
environments~\cite{wang2025openhands}. Terminal-Bench extends executable
evaluation to diverse command-line workflows, with task-specific environments and
tests~\cite{merrill2026terminalbench}; its scope is broader than repository repair.
Harness-Bench studies variation across model--harness configurations under shared
task environments and budgets~\cite{yao2026harnessbench}.

These studies motivate treating the execution configuration as part of the
measured system. We use a fixed OpenHands protocol for the main comparison and
interpret each result as a model--harness outcome. The contribution is a task
definition, curated benchmark, and analysis of delivered artifacts under that
protocol. It is not a new agent runtime, and the reported model ordering is not
assumed to transfer unchanged to other tool interfaces or execution policies.
